\section{Method}

12-lead ECG data was first obtained from the Physikalisch-Technische Bundesanstalt (PTB) Diagnostic ECG Database [FIXME].
Each record was sampled at 1 kHz, and a duration of 1 second (\(N=1000\)) was selected for the sake of brevity.
\newline
\par
To produce the compressed representation of the ECG, the sensing matrix \(\Phi\) was constructed.
The exact details of this matrix varies by implementation, but this paper follows the implementation described by Daponte et al. [FIXME].
Here, the authors propose that \(\Phi\) be composed of a combination of dominant leads.
Although less practical, this faired notably better than alternate methods, such an iid Gaussian sensing matrix [FIXME].
For the given ECG frame, a reference was constructed by taking the root mean square of leads II and aVF, \(x_0\), and then centering the signal subject to the following equations.

\[ x_0 = \sqrt{\frac{1}{2}(x_{II}^2 + x_{aVF}^2)} \]
\[ x_p = | x_0 - \bar{x}_0 | \]

A binary mask was applied to \(x_p\) with a threshold of 60\% to produce the vector \(p\).
This vector was then used to construct \(\Phi\) as follows:
\[
\Phi = \begin{bmatrix}
p(1) & p(2) & \cdots & p(N) \\
p(N-CR+1) & p(N-CR+2) & \cdots & p(N-CR) \\
\vdots & \vdots & \ddots & \vdots \\
p(CR+1) & p(CR+2) & \cdots & p(CR)
\end{bmatrix}.
\]
\newline
\noindent
The Python code for constructing the sensing matrix is shown below.

\begin{lstlisting}[language=Python, caption=Sensing Matrix Construction]
def sensing_matrix(ecg: Ecg, compression_ratio: int):
    # Construct a reference signal from Leads II, aVF
    x0 = numpy.sqrt(0.5 * (ecg.get('ii') ** 2 + ecg.get('avf') ** 2))

    # Center the reference signal and apply thresholding
    xp = numpy.abs(x0 - numpy.mean(x0))
    p = (xp >= numpy.percentile(xp, 60)).astype(int).flatten()

    # Construct the sensing matrix
    compressed_len = ecg.length // compression_ratio
    phi = numpy.zeros((compressed_len, ecg.length), dtype=int)
    for i in range(compressed_len):
        phi[i, :] = numpy.roll(p, -i * compression_ratio)
    return phi
\end{lstlisting}

\par
Applying the compression is now trivial, as the compressed data matrix \(Y\) is simply the product of the sensing matrix and ECG data.
\[ Y = \Phi X \]

At the receiver, recovery of the original ECG source can begin.
First, a wavelet basis \(\Psi\) must be constructed such that
\[ X = \Psi C \]
where \(C\) is a sparse matrix of coefficients.
Again, this paper follows the implementation described by Daponte et al. [FIXME], selecting the following basis matrix
\[ \Psi = [\Psi', u]\]
where \(u\) is an \(N\)-length column \(\frac{1}{\sqrt{N}} \cdot [1,... ,1]^T\) and \(\Psi'\) contains columns of mexican hat wavelets in the form:
\[ \psi(a,b) = \frac{2}{\sqrt{3a}\,\pi^{1/4}} \cdot \left[1 - \left(\frac{n - b}{a}\right)^2 \right] \cdot e^{-\frac{1}{2}\left(\frac{n - b}{a}\right)^2}, \]
\noindent
The Python code for constructing the wavelet basis matrix is shown below.

\begin{lstlisting}[language=Python, caption=Wavelet Basis Matrix Construction]
def mexican_hat(length: int, a: float, b: float):
    n = numpy.arange(length)
    # Construct wavelet
    scalar = (2 / (numpy.sqrt(3 * a) * (numpy.pi ** 0.25)))
    n_min_b_ovr_a_sq = ((n - b) / a) ** 2
    wavelet = scalar * (1 - n_min_b_ovr_a_sq) * numpy.exp(-0.5 * n_min_b_ovr_a_sq)
    return wavelet

def basis_matrix(length: int):
    cols = []
    # Construt wavelet dictionary
    for m in range(1, int(numpy.floor(numpy.log2(length))) + 1):
        a = 2 ** m
        for k in range((length - 1) // a + 1):
            b = k * a
            psi = mexican_hat(length, a, b)
            cols.append(psi)
    # Pad with an additional column
    u = (1 / numpy.sqrt(length)) * numpy.ones(length)
    return numpy.column_stack(cols)
\end{lstlisting}

Given the above information, a sparse minimization problem can now be designed to solve for the coefficient matrix \(C\) with cost function \(J\) [FIXME]:
\[ \hat{C} = \arg\min_{C}\left\|Y - \Phi \Psi C\right\|_F^2 + \lambda J \]
\noindent
To demonstrate the trade-offs and effectiveness of each algorithm for recovering the ECG signal, the following cost functions are investigated:
\newline
\indent \(\ell_1\)-norm (LASSO) [FIXME]:
\[ J = \left\| C \right\|_1 \]
\indent \(\ell_{2,1}\)-norm (Group LASSO) [FIXME]:
\[ J = \left\| C \right\|_2,1 = \sum_{n=1}^N \left\| c_{n\ \cdot} \right\|_2 \]
\indent \(\ell_{2,p}\)-norm (M-FOCUSS) [FIXME]:
\[ J = \sum_{n=1}^N \left\| c_{n\ \cdot} \right\|_2^p,\ p \in [0,1] \]
\noindent
The Python code for each implementation is shown below:

\begin{lstlisting}[language=Python, caption=Implementations for Reconstruction Algorithms]
# LASSO / Group LASSO are convex optimization problems, which can be solved with the cvxpy library.

def lasso_solver(compressed_data, phi, psi, regularization):
    c_var = cvxpy.Variable((psi.shape[1], compressed_data.shape[1]))
    cost_func = regularization * cvxpy.norm(c_var, 1)
    prob = cvxpy.Problem(
        cvxpy.Minimize(
            cvxpy.norm(compressed_data - (phi @ psi @ c_var), "fro") ** 2 + cost_func
        )
    )
    prob.solve(solver=cvxpy.SCS)
    return c_var.value

def group_lasso_solver(compressed_data, phi, psi, regularization):
    c_var = cvxpy.Variable((psi.shape[1], compressed_data.shape[1]))
    cost_func = regularization * cvxpy.sum(cvxpy.norm(c_var, 2, axis=1))
    prob = cvxpy.Problem(
        cvxpy.Minimize(
            cvxpy.norm(compressed_data - (phi @ psi @ c_var), "fro") ** 2 + cost_func
        )
    )
    prob.solve(solver=cvxpy.SCS)
    return c_var.value
\end{lstlisting}

Finally, the uncompressed ECG can be reconstructed by multiplying the wavelet basis with the recovered coefficient matrix:
\[ \hat{X} = \Psi \hat{C} \]
